\subsection{EXTENSION OF THE BASE RING}

Let \(\mathbf{K}_1\) be a commutative ring with unit element and \(\phi\) a homomorphism of K into \(\mathbf{K}_1\) mapping 1 to 1. Let \(\mathfrak{g}\) be a Lie K-algebra, U its enveloping algebra and M a left g-module, that is a left U-module. Then \(\mathrm{M}_{(\mathbf{K}_1)}\) has a canonical left \(\mathrm{U}_{(\mathbf{K}_1)}\) -module structure and hence a left \(\mathfrak{g}_{(\mathbf{K}_1)}\) -module structure. Let \(\rho\), \(\rho_{(\mathbf{K}_1)}\) be the representations of \(\mathfrak{g}\) and \(\mathfrak{g}_{(\mathbf{K}_1)}\) corresponding to M and \(\mathrm{M}_{(\mathbf{K}_1)}\) ; \(\rho_{(\mathbf{K}_1)}\) is said to be derived from \(\rho\) by extending the base ring and the results of Algebra, Chapter VIII, § 13, no. 4 can be applied. If \(x \in \mathfrak{g}\) , \(\rho_{(\mathbf{K}_1)}(x)\) is just the endomorphism \(\rho(x) \otimes 1\) of \(\mathrm{M}_{(\mathbf{K}_1)} = \mathrm{M} \otimes_{\mathbf{K}} \mathbf{K}_1\) .

Suppose that K is a field, that \(K_{1}\) is an extension of K and that \(\phi\) is the canonical injection of K into \(K_{1}\) . Let V and \(V'\) be vector subspaces of M. Let a be the vector subspace of g consisting of the \(x \in g\) such that \(\rho(x)(V) \subset V'\) . Let \(a'\) be the vector subspace of \(\mathfrak{g}_{(\mathbf{K}_{1})}\) consisting of the \(x' \in \mathfrak{g}_{(\mathbf{K}_{1})}\) such that \(\rho_{(\mathbf{K}_{1})}(x')(V_{(\mathbf{K}_{1})}) \subset V'_{(\mathbf{K}_{1})}\) . Then \(a' = a_{(\mathbf{K}_{1})}\) . For clearly \(a_{(\mathbf{K}_{1})} \subset a'\) . Now let \(x' \in a'\) . We may write \(x' = \sum_{i=1}^{n} \lambda_{i} x_{i}\) , where the \(x_{i}\) are in g and the \(\lambda_{i}\) are elements of \(K_{1}\) linearly independent over K. For all \(u \in V\) , \(\rho(x') .u \in V'_{(\mathbf{K}_{1})}\) , that is \(\sum_{i=1}^{n} \lambda_{i} \rho(x_{i}) .u \in V'_{(\mathbf{K}_{1})}\) , whence \(\rho(x_{i}) .u \in V'\) , hence \(x_{i} \in a\) and \(x' \in a_{(\mathbf{K}_{1})}\) . This shows that \(a' = a_{(\mathbf{K}_{1})}\) . In particular, the centre of \(\mathfrak{g}_{(\mathbf{K}_{1})}\) is derived from the centre of g by extending K to \(K_{1}\) : it suffices to apply the above to the adjoint representation of g. It follows that \(\mathcal{C}_{\rho'}(\mathfrak{g}_{(\mathbf{K}_{1})}) = (\mathcal{C}_{\rho}\mathfrak{g})_{(\mathbf{K}_{1})}\) for all p.~Similarly, let h be a subalgebra of g and n the normalizer of h in g. Then the normalizer of \(\mathfrak{h}_{(\mathbf{K}_{1})}\) in \(\mathfrak{g}_{(\mathbf{K}_{1})}\) is \(n_{(\mathbf{K}_{1})}\) .

Let \(\mathbf{K}, \mathbf{K}_{1}, \mathfrak{g}, \rho, \mathbf{M}\) be as in the last paragraph. Let \(\mathfrak{b}\) be a vector subspace of \(\mathfrak{g}\) and \(\mathbf{W}\) a vector subspace of \(\mathbf{M}\) . Let \(\mathbf{V}\) be the vector subspace of \(\mathbf{M}\) consisting of the \(m \in \mathbf{M}\) such that \(\rho(\mathfrak{b}). m \subset W\) . Let \(V'\) be the vector subspace of \(\mathrm{M}_{(\mathbf{K}_{1})}\) consisting of the \(m' \in \mathrm{M}_{(\mathbf{K}_{1})}\) such that \(\rho_{(\mathbf{K}_{1})}(\mathfrak{b}_{(\mathbf{K}_{1})}). m' \subset W_{(\mathbf{K}_{1})}\) . As above it is seen that \(V' = V_{(K_1)}\) . In particular, the vector subspace of invariants of \(M_{(K_1)}\) is derived from the vector subspace of invariants of M by extending the base field from K to \(K_1\) .

Let \(K, K_{1}\) and \(\phi\) be as at the beginning of this no. Let g be a Lie K-algebra and M and N g-modules. If M and N are isomorphic g-modules, \(M_{(K_{1})}\) and \(N_{(K_{1})}\) are isomorphic \(g_{(K_{1})}\) -modules. Conversely:

PROPOSITION 13. Let K be a field, \(K_{1}\) an extension of K, g a Lie K-algebra and M, N two g-modules of finite dimension over K. If \(\mathrm{M}_{(\mathbf{K}_{1})}\) and \(\mathrm{N}_{(\mathbf{K}_{1})}\) are isomorphic \(\mathfrak{g}_{(\mathbf{K}_{1})}\) -modules, M and N are isomorphic g-modules.

The proof is in two steps.

\begin{enumerate}
\def\labelenumi{(\arabic{enumi})}
\item
  Suppose first that \(\mathbf{K}_1\) is an extension of K of finite degree \(n\) . Let U be the enveloping algebra of \(\mathfrak{g}\) , so that the enveloping algebra of \(\mathfrak{g}_{(\mathbf{K}_1)}\) is \(\mathrm{U}_{(\mathbf{K}_1)} = \mathrm{U} \otimes_{\mathbb{K}} \mathrm{K}_1\) (§ 2, no. 9). As \(\mathrm{M}_{(\mathbf{K}_1)}\) and \(\mathrm{N}_{(\mathbf{K}_1)}\) are isomorphic as \(\mathrm{U}_{(\mathbf{K}_1)}\) -modules they are a fortiori isomorphic as U-modules; but as U-modules they are respectively isomorphic to \(\mathbf{M}^n\) and \(\mathbf{N}^n\) . Now M and N are U-modules of finite length; M (resp. N) is therefore the direct sum of a family \((\mathrm{P}_i^{s_i})_{1 \leqslant i \leqslant p}\) (resp. \((\mathrm{Q}_j^{s_j})_{1 \leqslant j \leqslant q}\) ) of submodules such that the \(\mathrm{P}_i\) (resp. \(\mathrm{Q}_j\) ) are indecomposable and two \(\mathrm{P}_i\) (resp. \(\mathrm{Q}_j\) ) of different indices are not isomorphic (Algebra, Chapter VIII, § 2, no. 2, Theorem 1). Then \(\mathrm{M}^n\) (resp. \(\mathrm{N}^n\) ) is isomorphic to the direct sum of the \(\mathrm{P}_i^{nr_i}\) (resp. \(\mathrm{Q}_j^{ns_i}\) ); it follows (loc. cit.) that \(p = q\) and that after permuting the \(\mathrm{Q}_j\) if necessary \(nr_i = ns_i\) and \(\mathrm{P}_i\) is isomorphic to \(\mathrm{Q}_i\) for \(1 \leqslant i \leqslant p\) , hence M is isomorphic to N.
\item
  General case. Let P be the g-module \(\mathcal{L}_{\mathrm{K}}(\mathrm{M},\mathrm{N})\) and Q the subspace of invariants of P, that is the set of homomorphisms of the g-module M into the g-module N. In the \(\mathfrak{g}_{(\mathrm{K}_{1})}\) -module \(\mathcal{L}_{\mathrm{K}_{1}}(\mathrm{M}_{(\mathrm{K}_{1})},\mathrm{N}_{(\mathrm{K}_{1})})=(\mathcal{L}_{\mathrm{K}}(\mathrm{M},\mathrm{N})_{\mathrm{K}_{1}})\) , the subspace of invariants is \(\mathrm{Q}_{(\mathrm{K}_{1})}\) . The hypothesis that \(\mathrm{M}_{(\mathrm{K}_{1})}\) and \(\mathrm{N}_{(\mathrm{K}_{1})}\) are isomorphic implies that M and N have the same dimension over K and that there exists in \(\mathrm{Q}_{(\mathrm{K}_{1})}\) an element g which is an isomorphism of \(\mathrm{M}_{(\mathrm{K}_{1})}\) onto \(\mathrm{N}_{(\mathrm{K}_{1})}\) . Let \((f_{1},\ldots,f_{d})\) be a basis of Q over K and choose bases of M and N over K. If
\end{enumerate}

\[
f = \sum_ {k = 1} ^ {d} \lambda_ {k} f _ {k}
\]

\(\lambda_{k}\in K_{1}\) for \(1\leqslant k\leqslant d\) , the matrix of \(f=\sum_{k=1}\lambda_{k}f_{k}\) with respect to these bases has determinant which is a polynomial \(D(\lambda_{1},\ldots,\lambda_{d})\) with coefficients in K. When f=g, this determinant is non-zero and hence the coefficients of D are not all non-zero. Therefore, if \(\Omega\) is the algebraic closure of K, there exists (since \(\Omega\) is infinite) elements \(\mu_{k}\in\Omega\) ( \(1\leqslant k\leqslant d\) ) such that \(D(\mu_{1},\ldots,\mu_{d})\neq0\) (Algebra, Chapter IV, § 2, no. 5, Proposition 8). If \(K_{2}\) is the algebraic extension

\[
\sum_ {k = 1} ^ {d} \mu_ {k} f _ {k}
\]

of K generated by the \(\mu_{k}\) ( \(1 \leqslant k \leqslant d\) ), it follows that \(\sum_{k=1}^{d} \mu_{k} f_{k}\) is an isomorphism of \(\mathrm{M}_{(\mathbf{K}_{2})}\) onto \(\mathrm{N}_{(\mathbf{K}_{2})}\) ; but \(\mathbf{K}_{2}\) is of finite degree over K (Algebra, Chapter V, § 3, no. 2, Proposition 5) and hence M and N are isomorphic by the first part of the argument.

Again let \(K, K_{1}\) and \(\phi\) be as the beginning of this no. Let \(\rho\) be a representation of \(\mathfrak{g}\) on a K-module M with a finite basis \((x_{1},\ldots ,x_{n})\) . Then the bilinear form on \(\mathfrak{g}_{(\mathrm{K}_1)}\) associated with \(\rho_{(\mathrm{K}_1)}\) is derived from the bilinear form associated with \(\rho\) by extending the base ring to \(\mathrm{K}_1\) (for, if \(u\in \mathcal{L}_{\mathrm{K}}(\mathbf{M})\) , \(u\) has the same matrix with respect to \((x_{1},\dots,x_{n})\) as \(u\otimes 1\) with respect to \((x_{1}\otimes 1,\dots,x_{n}\otimes 1)\) and hence \(u\) and \(u\otimes 1\) have the same trace). In particular, if the K-module \(\mathfrak{g}\) has a finite basis, the Killing form of \(\mathfrak{g}_{(\mathrm{K}_1)}\) is derived from that of \(\mathfrak{g}\) by extending the base ring to \(\mathrm{K}_1\) .

